\begin{helper}
For every pair of natural numbers \(s,k\), the set
\(\{n\in\mathbb{N}:\mathrm{RamseyProperty}(s,k,n)\}\) is nonempty.
\end{helper}
\begin{proof}
We prove the claim by induction on \(s\), with an inner induction on \(k\) in
the successor case.

If \(s=0\), then \(n=0\) has \(\mathrm{RamseyProperty}(0,k,0)\) for every
\(k\): for any graph on \(\operatorname{Fin}(0)\), the empty finite set is a
\(0\)-clique.  This is exactly the first alternative in
\(\noderef{RamseyProperty}\).

Now suppose the claim has been proved for a fixed \(s\) and all \(k\).  We
prove it for \(s+1\) by induction on \(k\).  If \(k=0\), then \(n=0\) has
\(\mathrm{RamseyProperty}(s+1,0,0)\), because for any graph on
\(\operatorname{Fin}(0)\), the empty finite set is a \(0\)-clique in the
complement graph.  This is the second alternative in
\(\noderef{RamseyProperty}\).

For the successor step, assume the set of \(n\) satisfying
\(\mathrm{RamseyProperty}(s+1,k,n)\) is nonempty, and use the outer induction
hypothesis to choose \(a\) with \(\mathrm{RamseyProperty}(s,k+1,a)\).  Use the
inner induction hypothesis to choose \(b\) with
\(\mathrm{RamseyProperty}(s+1,k,b)\).  Then
\(\noderef{RamseyPropertyStep}\) gives
\(\mathrm{RamseyProperty}(s+1,k+1,a+b+1)\).  Hence the required set is
nonempty in the successor case as well.

By the two inductions, the set is nonempty for all natural numbers \(s,k\).
\end{proof}
